\section{The sporadic simple groups}
\label{sec:sporadic}

\subsection{Rational factors for the sporadic simple groups}

Names and class counts are those of the \emph{Atlas}~\cite{Atlas}.
Write
\[
 K_{-3}=\Q(\sqrt{-3}),\qquad K_5=\Q(\zeta_5).
\]
The degree of $K_5/\Q$ is four.  Hence a factor $M_r(K_5)$ gives four
conjugate copies of $M_r(\C)$ after scalar extension.

\begin{lemma}
\label{lem:schur-index}
Let $B=M_r(D)$ be a simple rational factor with center $K$, Schur index
$s=\sqrt{\dim_KD}$, and split degree $d=rs$.  Suppose the associated
$B$-module in rational class-coordinate space has complex multiplicity $m$.
Then $s$ divides both $d$ and $m$.  In particular, $\gcd(d,m)=1$ forces
$D=K$.  For the possible indices in the rational group-algebra setting
see~\cite{Yamada}.
\end{lemma}

\begin{proof}
For each embedding $K\hookrightarrow\C$, scalar extension sends $B$ to
$M_{rs}(\C)$.  A simple left $B$-module has $K$-dimension $rs^2$ and
complexifies to $s$ copies of the natural $rs$-dimensional module
\cite[\S{}29]{Reiner}.  Every
rational $B$-module therefore gives a complex multiplicity divisible by
$s$, while $s$ also divides $d=rs$.
\end{proof}

\begin{theorem}
\label{thm:sporadic-census}
The rational algebras of the twenty-six sporadic simple
groups are listed below.  The columns $z$ and $c$ give $\dim_{\Q}Z(A)$ and
$\dim_{\Q}A'$.  Every simple factor has Schur index one.  The only center
fields are $\Q$, $K_{-3}$, and $K_5$, and $K_5$ occurs only for $Ly$.

\begingroup
\footnotesize
\normalfont
\setlength{\tabcolsep}{3pt}
\renewcommand{\arraystretch}{1.08}
\begin{longtable}{@{}lrrrrp{0.54\textwidth}@{}}
\toprule
$G$ & $k(G)$ & $\dim A$ & $z$ & $c$ & $A$ over $\Q$\\
\midrule
\endfirsthead
\toprule
$G$ & $k(G)$ & $\dim A$ & $z$ & $c$ & $A$ over $\Q$\\
\midrule
\endhead
$M_{11}$ & 10 & 66 & 3 & 3 & $M_8(\Q)\oplus\Q^2$\\
$M_{12}$ & 15 & 149 & 3 & 3 & $M_{12}(\Q)\oplus M_2(\Q)\oplus\Q$\\
$M_{22}$ & 12 & 102 & 3 & 3 & $M_{10}(\Q)\oplus\Q^2$\\
$M_{23}$ & 17 & 151 & 5 & 5 & $M_{12}(\Q)\oplus M_2(\Q)\oplus\Q^3$\\
$M_{24}$ & 26 & 452 & 4 & 4 & $M_{21}(\Q)\oplus M_3(\Q)\oplus\Q^2$\\
$J_1$ & 15 & 111 & 4 & 4 & $M_{10}(\Q)\oplus M_3(\Q)\oplus K_{-3}$\\
$J_2$ & 21 & 269 & 3 & 3 & $M_{16}(\Q)\oplus M_3(\Q)\oplus M_2(\Q)$\\
$J_3$ & 21 & 209 & 6 & 6 & $M_{14}(\Q)\oplus M_3(\Q)\oplus\Q^2\oplus K_{-3}$\\
$J_4$ & 62 & 2000 & 11 & 11 &
$M_{44}(\Q)\oplus M_7(\Q)\oplus M_2(\Q)^2\oplus\Q\oplus K_{-3}^3$\\
$HS$ & 24 & 444 & 4 & 4 & $M_{21}(\Q)\oplus\Q^3$\\
$McL$ & 24 & 334 & 5 & 5 & $M_{18}(\Q)\oplus M_2(\Q)^2\oplus\Q^2$\\
$He$ & 33 & 647 & 5 & 5 & $M_{25}(\Q)\oplus M_4(\Q)\oplus M_2(\Q)\oplus\Q^2$\\
$Ru$ & 36 & 792 & 9 & 9 & $M_{28}(\Q)\oplus\Q^4\oplus K_{-3}^2$\\
$Suz$ & 43 & 1381 & 5 & 5 & $M_{37}(\Q)\oplus M_3(\Q)\oplus\Q^3$\\
$O'N$ & 30 & 411 & 9 & 12 & $M_{20}(\Q)\oplus M_2(\Q)\oplus\Q^5\oplus K_{-3}$\\
$Co_1$ & 101 & 9415 & 4 & 4 & $M_{97}(\Q)\oplus M_2(\Q)\oplus\Q^2$\\
$Co_2$ & 60 & 2926 & 5 & 5 & $M_{54}(\Q)\oplus M_2(\Q)^2\oplus\Q^2$\\
$Co_3$ & 42 & 1306 & 5 & 5 & $M_{36}(\Q)\oplus M_2(\Q)^2\oplus\Q^2$\\
$Fi_{22}$ & 65 & 3261 & 7 & 7 & $M_{57}(\Q)\oplus M_2(\Q)^2\oplus\Q^4$\\
$Fi_{23}$ & 98 & 8294 & 6 & 6 & $M_{91}(\Q)\oplus M_3(\Q)\oplus\Q^4$\\
$Fi_{24}'$ & 108 & 8853 & 12 & 15 & $M_{94}(\Q)\oplus M_2(\Q)^2\oplus\Q^9$\\
$HN$ & 54 & 1879 & 7 & 10 & $M_{43}(\Q)\oplus M_5(\Q)\oplus\Q^5$\\
$Ly$ & 53 & 1543 & 12 & 12 &
$M_{39}(\Q)\oplus M_3(\Q)\oplus M_2(\Q)\oplus\Q^3\oplus K_{-3}\oplus K_5$\\
$Th$ & 48 & 1384 & 10 & 10 & $M_{37}(\Q)\oplus M_2(\Q)^2\oplus\Q^7$\\
$\B$ & 184 & 26598 & 13 & 19 &
$M_{163}(\Q)\oplus M_3(\Q)\oplus M_2(\Q)^3\oplus\Q^8$\\
$\M$ & 194 & 28958 & 15 & 15 &
$M_{170}(\Q)\oplus M_5(\Q)\oplus M_3(\Q)\oplus M_2(\Q)^4\oplus\Q^8$\\
\bottomrule
\end{longtable}
\endgroup
\end{theorem}

\begin{proof}
The computation described in \cref{sec:certification} constructs each
rational word algebra and verifies closure.  A separating central element
gives square-free center factors; rational central idempotents give the
dimensions of the corresponding rational ideals.  These data determine the
split degree and class-space multiplicity of each complex block.  Their
gcd is one, so \cref{lem:schur-index} gives Schur
index one.

The quadratic defining polynomials have square-free discriminant $-3$.  The
only quartic center is defined by
$x^4+x^3+x^2+x+1$, which gives $K_5$.  Ranks at the primes
$1{,}000{,}003$ and $1{,}000{,}033$ reproduce the algebra, center, and
commutant dimensions of each residual calculation.
\end{proof}

\begin{corollary}
\label{cor:sporadic-saturation}
For each of the twenty-two sporadic groups whose class-coordinate
representation is multiplicity free, \(\Pram(G)=\PAd(G)\).  In particular the
Monster satisfies \(\Pram(\M)=\PAd(\M)\).
\end{corollary}

\begin{proof}
Immediate from \cref{thm:good-label-saturation} and the equality \(z=c\)
recorded in the table.
\end{proof}


The class-coordinate representation is multiplicity-free for 22 sporadic
groups.  The exceptions are $O'N$, $Fi_{24}'$, $HN$, and $\B$.
For the first three, a scalar complex block occurs with multiplicity two.
For $\B$, one scalar block and one $3$-dimensional block each occur
with multiplicity two.  The distinction between $z$ and $c$ in the table
makes these multiplicities visible.

\subsection{Good-label characters}
\label{sec:characters}

The sporadic census contains nine blocks with center \(\Q(\zeta_3)\) and
one block with center \(\Q(\zeta_5)\).

Throughout, \(N=\exp(G)\), \(U_N=(\Z/N)^\times\), \(\Gamma=\Psi(U_N)\), and
\(\chi_D\) denotes the Kronecker symbol \(\bigl(\tfrac D\cdot\bigr)\) of a
fundamental discriminant \(D\).

\subsubsection{Recovering the unit action from ramified data}

The CTblLib input consists of class sizes, element orders, and the power
maps \(\sigma_p\) for primes \(p\mid N\).

\begin{definition}
\label{def:witness-subgroup}
For an integer \(m\mid N\) put
\begin{equation}
 H_m=\bigl\langle\,p\bmod m\ :\ p\ \text{prime},\ p\mid N,\ p\nmid m\,\bigr\rangle
 \ \leq\ U_m .
 \label{eq:witness-subgroup}
\end{equation}
We call \(G\) \emph{ramified transparent} if \(H_m=U_m\) for every element
order \(m\) of \(G\).
\end{definition}

\begin{theorem}
\label{thm:unit-reconstruction}
Let \(C\) be a conjugacy class whose elements have order \(m\), let
\(u\in U_N\), and let \(v\) be any ramified label with \(v\equiv u\pmod m\).
Then
\begin{equation}
 \sigma_u(C)=\sigma_v(C).
 \label{eq:unit-reconstruction}
\end{equation}
If \(G\) is ramified transparent, the maps
\(\{\sigma_p:p\mid N\}\) determine the action of \(U_N\) on \(\Conj(G)\).
\end{theorem}

\begin{proof}
For \(x\in C\) we have \(x^m=1\), so \(x^u=x^{u\bmod m}=x^{v\bmod m}=x^v\),
whence \(\sigma_u(C)=\sigma_v(C)\), which is
\eqref{eq:unit-reconstruction}.  A ramified \(v\) with \(v\equiv u\pmod m\)
exists if and only if \(u\bmod m\in H_m\), because the residues of ramified
labels modulo \(m\) that are units are the elements of \(H_m\):
a ramified \(v\) is a product of primes dividing \(N\), those dividing \(m\)
contribute non-units, and \(\gcd(v,m)=1\) forces every prime factor of \(v\)
to be coprime to \(m\).  Under ramified transparency this holds for all
\(m\) and \(u\), and \(\sigma_v\) is computed from the stored maps by
multiplicativity, \(\sigma_v=\prod_p\sigma_p^{e_p}\) for \(v=\prod_pp^{e_p}\).
\end{proof}

\begin{proposition}
\label{prop:reconstruction-coverage}
Twenty-four of the twenty-six sporadic simple groups are ramified
transparent.  The exceptions are \(J_3\), where
\([U_{15}:H_{15}]=2\), and \(Th\), where \([U_{30}:H_{30}]=2\); in both cases
the missing coset is the one containing \(-1\), and the ambiguity concerns a
single pair of classes.
\end{proposition}

\begin{proof}
Closing the residues of the primes dividing \(N\) under multiplication
modulo \(m\) and comparing \(|H_m|\) with \(\varphi(m)\) leaves only the two
displayed failures.  For \(J_3\), \(N=2^3\cdot3^2\cdot
5\cdot17\cdot19\) and the primes coprime to \(15\) reduce to
\(2,17,19\equiv2,2,4\pmod{15}\), generating
\(\langle2\rangle=\{1,2,4,8\}\) of index two in \(U_{15}\); the missing coset
is \(7\langle2\rangle=\{7,11,13,14\}\ni-1\).  The computation for \(Th\) at
\(m=30\) is identical in form.
\end{proof}

For \(J_3\) and \(Th\), call a permutation of \(\Conj(G)\)
\emph{admissible} if it extends the forced values, preserves class size and
element order, and commutes with the stored power maps.

\begin{proposition}
\label{prop:ambiguity-harmless}
For each sporadic simple group and each chosen generator of \(U_N\), the
admissible completion is unique.  The stored ramified maps, class sizes, and
element orders determine the action of \(U_N\) on \(\Conj(G)\).
\end{proposition}

\begin{proof}
For the twenty-four ramified transparent groups,
\cref{thm:unit-reconstruction} determines all images.  For \(J_3\), two of
the six chosen generators, those with nontrivial \(3\)- and
\(5\)-components, leave undetermined the two
classes of element order \(15\), the remaining nineteen images being forced;
of the two possible assignments only one commutes with \(\sigma_3\).  For
\(Th\), one of the eight generators, the \(3\)-component, leaves the
two classes of element order \(30\) undetermined out of forty-eight, and
there the disambiguating constraint is \(\sigma_2\).  Which of the two
assignments survives is not uniform: the surviving completion fixes both
classes for the \(3\)-component of \(J_3\) and transposes them in the other
two cases.
\end{proof}

The statement is independent of the displayed generators: once the
permutations are recovered on any generating set, the homomorphism
\(U_N\to\operatorname{Sym}(\Conj(G))\) determines every other choice.

\subsubsection{The character census}

\begin{theorem}
\label{thm:good-label-census}
Let \(G\) be a sporadic simple group and \(E_k\) a primitive central
idempotent of \(\Pram(G)\).  There is a character
\(\lambda_k:U_N\to\mu(K_k)\) with
\[
 \Psi_uE_k=\lambda_k(u)E_k .
\]
The principal block carries the trivial character.  Of the remaining blocks,
ten have nonrational center and are listed in
\cref{sec:nonrational-blocks}; every other one carries the trivial character
or the Kronecker symbol \(\chi_{D_k}\) of a fundamental discriminant \(D_k\)
whose prime divisors divide \(N\).  Five blocks, in \(O'N\), \(Fi_{24}'\),
\(HN\), and \(\B\), have multiplicity two, and \(\Psi_u\) is scalar on those
as well.
\end{theorem}

\begin{proof}
Compute the primitive central idempotents over \(\Q\) and the operators
\(\Psi_u\) for one generator per cyclic factor of
\(\prod_{p^a\|N}U_{p^a}\).  The block-local theorem treats the
multiplicity-one blocks.  Rational row reduction gives
\(\Psi_uE_k=\lambda_k(u)E_k\) on the five multiplicity-two blocks.  The
order of \(\lambda_k(u)\) is the least \(e\) with
\((\Psi_uE_k)^e=E_k\); the character
order is the least common multiple over the generators; the conductor is
assembled from the local orders by the standard rule for
\(U_{p^a}\).  For a rational block the parity \(\lambda_k(-1)\) selects the
sign of the fundamental discriminant of the given conductor.  Direct
substitution gives \(\lambda_k(u)=\chi_D(u)\) on the generators.

The principal block character is trivial.  That block is the cyclic
module \(A\delta_{1A}\) generated by the identity-class coordinate
(\cref{sec:certification}); every \(\Psi_u\) fixes \(\delta_{1A}\) and
commutes with \(A\), so it fixes \(A\delta_{1A}\) pointwise.
\end{proof}

\subsubsection{The ten nonrational blocks}
\label{sec:nonrational-blocks}

The ten nonrational blocks occur in
\[
 J_1,\qquad J_3,\qquad J_4,\qquad Ru,\qquad O'N,\qquad Ly.
\]
Nine blocks have center \(\Q(\zeta_3)\); the unique remaining block has
center \(\Q(\zeta_5)\) and belongs to \(Ly\).  In every case the star acts as
complex conjugation on the center, the character order is the odd prime
\(e\) with \(K=\Q(\zeta_e)\), and the block is supported on \(e\)
classes of a single element order.  Its split degree is
\(d_k=\varphi(e)=e-1\).  The complete list is:

\begin{center}
\normalfont
\begin{tabular}{@{}llrrr@{}}
\toprule
$G$ & $K_k$ & $e$ & cond. & elt.\ order\\
\midrule
$J_1$  & $K_{-3}$ & 3 & 19 & 19\\
$J_3$  & $K_{-3}$ & 3 &  9 &  9\\
$J_4$  & $K_{-3}$ & 3 & 37 & 37\\
$J_4$  & $K_{-3}$ & 3 & 31 & 31\\
$J_4$  & $K_{-3}$ & 3 & 43 & 43\\
$Ru$   & $K_{-3}$ & 3 &  7 & 14\\
$Ru$   & $K_{-3}$ & 3 & 13 & 26\\
$O'N$  & $K_{-3}$ & 3 & 19 & 19\\
$Ly$   & $K_{-3}$ & 3 & 67 & 67\\
$Ly$   & $K_5$    & 5 & 31 & 31\\
\bottomrule
\end{tabular}
\end{center}


\begin{remark}
\label{rem:pariahs}
The six groups carrying a nonrational center are exactly the six pariah
groups, the sporadic groups not involved in the Monster.  Every group of
the happy family has \(\Pram(G)\) split over \(\Q\) with all centers
rational.  We record this as an observation of the census; no explanation
is offered here.
\end{remark}


\begin{proposition}
\label{prop:conductor-determined}
Let \(E\) be a block whose support consists of classes of a common element
order \(m\), and let \(\lambda_E\) be its good-label character, of order
\(e>1\).  Write \(m=2^am'\) with \(m'\) odd.  Then
\begin{enumerate}
\item \(\lambda_E\) factors through \(U_m\), so its conductor divides \(m\);
\item if \(a\leq1\), the conductor divides \(m'\);
\item if moreover \(m'=p^b\) is a prime power and
\(e\nmid\varphi(p^{b-1})\), the conductor equals \(m'\).
\end{enumerate}
\end{proposition}

\begin{proof}
(1) By \cref{thm:unit-reconstruction} the permutation \(\sigma_u\) acts on a
class of element order \(m\) through \(u\bmod m\) alone, so \(\lambda_E\)
factors through \(U_m\) and its conductor divides \(m\).

(2) For odd \(m'\) the Chinese remainder theorem gives
\(U_{2m'}\cong U_2\times U_{m'}=U_{m'}\), so when \(a\leq1\) the character
factors through \(U_{m'}\).

(3) The conductor is \(p^c\) for some \(c\leq b\), minimal such that
\(\lambda_E\) factors through \(U_{p^c}\); then
\(e=|\lambda_E(U_N)|\) divides \(|U_{p^c}|=\varphi(p^c)\).  For \(p\) odd and
\(c\leq b-1\) one has \(\varphi(p^c)\mid\varphi(p^{b-1})\), whence
\(e\mid\varphi(p^{b-1})\), contrary to hypothesis.  Hence \(c=b\).
\end{proof}

\begin{corollary}
\label{cor:conductor-odd-part}
In every row of \cref{sec:nonrational-blocks} the conductor equals the odd
part \(m'\) of the common element order.
\end{corollary}

\begin{proof}
Each row has \(a\leq1\), and \(m'\) is a prime power.  In nine rows \(m'\) is
one of \(7,13,19,31,37,43,67\); then \(b=1\), and \(e>1\) does not divide
\(\varphi(p^0)=1\).  In the remaining row \(m'=9\) and \(e=3\), which does
not divide \(\varphi(3)=2\).
\Cref{prop:conductor-determined} applies in every case.
\end{proof}

\begin{example}
\label{ex:ly-quintic}
Among the twenty-six sporadic groups, \(Ly\) alone has a quartic center.
Its five conjugacy classes of elements of order \(31\) form one rational
class and carry the full \(\Q(\zeta_5)\) factor of \(\Pram(Ly)\).  Writing
\(R\) for this five-class orbit, the block is
\[
 V_k=\Bigl\{f\in\Span_\Q(\delta_C:C\in R)\ :\ \textstyle\sum_{C\in R}f(C)=0\Bigr\},
 \qquad \dim_\Q V_k=4,
\]
and the good labels act through a quintic character of conductor \(31\).
Thus \(V_k\) is one-dimensional over its cyclotomic endomorphism field:
\[
 \End_{\Q[U_N]}(V_k)=\Q(\zeta_5).
\]
The character is even, so inversion fixes all five order-\(31\) classes.
\end{example}

\begin{example}
\label{ex:j1-cubic}
The three classes of elements of order \(19\) in \(J_1\) form a single
rational class; the augmentation-zero part of their span is a two-dimensional
block with center \(\Q(\sqrt{-3})=\Q(\zeta_3)\), on which the good labels act
through a cubic character of conductor \(19\).  Since \(19\equiv1\pmod3\),
such characters exist modulo \(19\); there are two of them, complex conjugate
to one another, and they are interchanged by the choice of embedding
\(\Q(\zeta_3)\hookrightarrow\C\).  The pair, not either member, is the
invariant.
\end{example}

\begin{theorem}
\label{thm:regular-rational-class}
Let \(R\subseteq\Conj(G)\) be a \(\Gamma\)-orbit with \(|R|=r\), let
\[
 W_R=\Bigl\{f\in\Span_\Q(\delta_C:C\in R):\textstyle\sum_{C\in R}f(C)=0\Bigr\},
 \qquad \dim_\Q W_R=r-1 .
\]
Then \(\Gamma\) acts regularly on \(R\) through an abelian quotient
\(A_R\) of order \(r\), and \(W_R\) is \(\Gamma\)-stable.  If \(r\) is
prime, then \(A_R\cong\Z/r\), and \(\Gamma\) acts on \(W_R\) through a
character of order \(r\) with values in \(\Q(\zeta_r)\), and
\(\End_{\Q[\Gamma]}(W_R)=\Q(\zeta_r)\).  If moreover \(W_R\) is a single
split block of \(\Pram(G)\) of multiplicity one, that block has center
\(K=\Q(\zeta_r)\) and good-label character of order \(r\).
\end{theorem}

\begin{proof}
\(\Gamma\) is abelian by \cref{lem:good-labels}(4) and acts transitively on
\(R\) by definition of an orbit.  In a transitive action of an abelian group
all point stabilizers agree; after quotienting by their common value
\(\Gamma_R\), the action is free and transitive.  Thus
\(A_R=\Gamma/\Gamma_R\) has order \(r\), and \(R\) is an \(A_R\)-torsor.
If \(r\) is prime then \(A_R\cong\Z/r\).  The span of \(R\) is the regular
representation \(\Q[A_R]\), which splits as \(\Q\oplus W_R\) with \(\Q\) the
invariants; both summands are \(\Gamma\)-stable.

Suppose \(r\) is prime.  Then \(\Q[\Z/r]\cong\Q\times\Q(\zeta_r)\)~\cite{Washington}, the
second factor acting on \(W_R\), so \(W_R\) is a one-dimensional
\(\Q(\zeta_r)\)-vector space on which a generator of \(\Z/r\) acts as
multiplication by \(\zeta_r\).  Hence the induced character has order
\(r\), its values generate \(\Q(\zeta_r)\), and the commuting algebra is
\(\Q(\zeta_r)\).  If \(W_R\) is a single split multiplicity-one block of
\(\Pram(G)\), its center is its commuting algebra inside the isotypic
component, which contains and is contained in \(\Q(\zeta_r)\) by the previous
sentence and \cref{thm:block-local-character}.
\end{proof}

\begin{corollary}
\label{cor:orbit-divisibility}
Let \(r\) be an odd prime.  Then \(V_G\) has a \(\Q[U_N]\)-summand with
endomorphism field \(\Q(\zeta_r)\) if and only if \(r\) divides \(|R|\) for
some rational class \(R\).
\end{corollary}

\begin{proof}
By \cref{thm:rational-class-decomposition},
\(V_G\cong\bigoplus_R\Q[A_R]\) with \(A_R=U_N/\Gamma_R\) abelian of order
\(|R|\).  For an abelian group \(A\), the simple summands of \(\Q[A]\) are the
fields \(\Q(\chi)=\Q(\zeta_{\mathrm{ord}\chi})\) as \(\chi\) runs over the
Galois orbits of characters of \(A\), each acting on itself, so the
endomorphism field of that summand is \(\Q(\chi)\).  Since \(r\) is an odd
prime, \(\Q(\chi)=\Q(\zeta_r)\) holds if and only if
\(\mathrm{ord}\,\chi\) is
\(r\) or \(2r\), and such a \(\chi\) exists for some \(A_R\) if and only if
\(r\) divides the exponent of some \(A_R\), equivalently \(r\mid|A_R|=|R|\).
\end{proof}

\begin{corollary}
\label{cor:nonrational-from-orbits}
Every nonrational block in the sporadic census arises as in
\cref{thm:regular-rational-class}: nine from rational classes of size three,
giving \(K=\Q(\zeta_3)=\Q(\sqrt{-3})\) and cubic characters, and one, the
\(\Q(\zeta_5)\) block of \(Ly\), from a rational class of size five, giving a
quintic character.  In every case \(d_k=r-1\).
\end{corollary}

\begin{proof}
For each block, compute its class support.  For all ten nonrational blocks the
support is a single orbit \(R\) of prime size \(r\in\{3,5\}\), all of whose
classes have one element order, and \(d_k=|R|-1\); the block therefore is
\(W_R\) and \cref{thm:regular-rational-class} applies.  Comparing its
conclusion with the computed columns, \(K_k=\Q(\zeta_r)\) and the character
order is \(r\), as tabulated.
\end{proof}

\Cref{thm:regular-rational-class} with \(r=2\) accounts for a large part of
the rational half of \cref{thm:good-label-census} as well: a rational class
splitting into two conjugacy classes contributes a one-dimensional block on
which \(\Gamma\) acts by a quadratic character.  The blocks with
\(d_k>r-1\) are the ones where several rational classes fuse into a single
simple factor; those blocks have a
conductor divisible by more than one prime.

\subsubsection{The Monster}

\begin{proposition}
\label{prop:monster-imaginary}
Let \(\M\) be the Monster, so that \(k(\M)=194\) and
\begin{equation}
 N=\exp(\M)=2^5\,3^3\,5^2\cdot7\cdot11\cdot13\cdot17\cdot19\cdot23
 \cdot29\cdot31\cdot41\cdot47\cdot59\cdot71 .
 \label{eq:monster-exponent}
\end{equation}
Its ramified algebra is
\[
 \Pram(\M)\cong
 M_{170}(\Q)\oplus M_5(\Q)\oplus M_3(\Q)
 \oplus M_2(\Q)^4\oplus\Q^8.
\]
Thus all fifteen blocks have center \(\Q\).  Two carry the trivial character;
the other thirteen carry quadratic characters whose
fundamental discriminants are
\begin{equation}
 -11,\,-20,\,-23,\,-31,\,-39,\,-47,\,-56,\,-59,\,-71,\,-87,\,-95,\,-104,\,-119 .
 \label{eq:monster-discriminants}
\end{equation}
All thirteen are negative, so every nontrivial good-label character of the
Monster is odd: \(\lambda_k(-1)=-1\) on each of the thirteen blocks.
\end{proposition}

\begin{proof}
The rational decomposition and the character conductors are the Monster row
of \cref{thm:sporadic-census,thm:good-label-census}.  The inversion map in
the CTblLib class ordering fixes \(172\) of the \(194\) classes, while the
unit-group action has \(172\) orbits.  Hence the numbers of real and rational
classes agree.  The defining module is multiplicity free, so
\cref{prop:real-rational-parity} forces every nontrivial block character to
be odd.  The thirteen fundamental discriminants are then determined by
their certified conductors and their values on the unit-group generators.
\end{proof}

By \cref{prop:rational-splitting}, the
\(+1\) eigenspace of \(\Psi_{-1}\) is spanned by the orbit sums of the
inversion action, so \eqref{eq:monster-discriminants} says that
every nontrivially twisted block of the Monster is detected already by
inversion: the twenty-two dimensions \(194-172\) lie entirely in the
\(-1\) eigenspace of complex conjugation on classes.  This is special to \(\mathbb M\) among the large
sporadic groups: \(Ly\), for instance, has four blocks with positive
discriminant \(40,37,24,21\) and one with \(-11\), and \(J_4\) has three
positive and one negative.

The assignment
\[
 G\ \longmapsto\ \bigl\{(K_k,\ d_k,\ \lambda_k)\bigr\}_{k}
\]
is an invariant of the marked weighted power system.  The groups \(Co_2\)
and \(Co_3\) have
the same rational block list \(M_r(\Q)\oplus M_2(\Q)^2\oplus\Q^2\) up to the
principal degree, but their residual characters differ, \(\chi_{-23},\chi_{-7}\)
and \(\chi_{-20},\chi_{-23}\), respectively.


\subsection{Rank-one saturation}
\label{sec:rank-one-saturation}

Let a finite elementary abelian \(2\)-group \(K\) act by weight-preserving
permutations on a finite weighted set \(X\).  Write
\[
 V_X=\bigoplus_{\chi\in\widehat K}V_\chi,
 \qquad
 V_\chi=\{v:kv=\chi(k)v\text{ for all }k\in K\}.
\]
Every \(V_\chi\) is rational.  Hence
\[
 \End_{\Q[K]}(V_X)=\bigoplus_{\chi:m_\chi>0}M_{m_\chi}(\Q),
 \qquad m_\chi=\dim_\Q V_\chi.
\]

\begin{lemma}[Cyclic rank-one generation]
\label{lem:cyclic-rank-one}
Let \(W\) be a finite-dimensional rational inner-product space, let
\(B\subseteq\End_\Q(W)\) be a unital star-algebra, and let \(v\in W\).
If the orthogonal rank-one projection \(E_v\) belongs to \(B\) and
\(Bv=W\), then \(B=\End_\Q(W)\).
\end{lemma}

\begin{proof}
Choose \(b_1,\ldots,b_d\in B\) such that \(b_1v,\ldots,b_dv\) is a basis
of \(W\).  The operators \(b_iE_vb_j^*\) are nonzero scalar multiples of
the matrix units for this basis.  They span \(\End_\Q(W)\).
\end{proof}

For \(a\) in the monoid generated by the labels, put
\[
 r_a=P_a^*P_a\mathbf1.
\]
The transfer identity makes \(P_a^*P_a=M_{r_a}\).  A finite tuple of the
functions \(r_a\) will be called a root profile.

\begin{lemma}[Root-profile projectors]
\label{lem:root-profile-projectors}
Let \(O\subseteq X\) be a \(K\)-orbit.  If a finite root profile takes one
joint value on \(O\) and a different joint value on every other orbit,
then the coordinate projection \(E_O\) belongs to \(\PE(\mathscr X)\).
\end{lemma}

\begin{proof}
The diagonal operators \(M_{r_a}\) belong to \(\PE(\mathscr X)\).
Multivariate Lagrange interpolation in the finitely many joint values
produces the indicator of \(O\).
\end{proof}

Suppose \(O=\{x,y\}\) is a two-point orbit belonging to the character
\(\chi\).  Since \(K\) preserves the weight, \(\mu(x)=\mu(y)\), and
\[
 E_O=E_{\delta_x+\delta_y}+E_{\delta_x-\delta_y}
\]
after normalizing the two orthogonal vectors.  Once the trivial matrix
block has been generated, an isolated root profile for \(O\) therefore
supplies the rank-one projection onto
\(\Q(\delta_x-\delta_y)\subseteq V_\chi\).

\begin{theorem}[Root-cyclic saturation]
\label{thm:root-cyclic-saturation}
Let \(A=\PE(\mathscr X)\subseteq\End_{\Q[K]}(V_X)\), where \(K\) is
elementary abelian of exponent \(2\).  Assume:
\begin{enumerate}
\item a fixed point \(x_0\) has \(E_{x_0}\in A\), and
      \(A\delta_{x_0}=V_{\mathbf1}\);
\item for each nontrivial \(\chi\) with \(m_\chi>1\), a two-point orbit
      \(O_\chi=\{x_\chi,y_\chi\}\) has an isolated root profile and
      \(A(\delta_{x_\chi}-\delta_{y_\chi})=V_\chi\);
\item the restrictions of elements of \(A\) separate the characters
      \(\chi\) with \(m_\chi=1\).
\end{enumerate}
Then
\[
 A=\End_{\Q[K]}(V_X).
\]
The saturation certificate requires
\[
 m_{\mathbf1}+
 \sum_{\substack{\chi\ne\mathbf1\\m_\chi>1}}m_\chi
\]
orbit vectors, together with scalar eigenvalue data on the
one-dimensional channels.
\end{theorem}

\begin{proof}
The first hypothesis and \cref{lem:cyclic-rank-one} generate the full
trivial block.  For every higher-dimensional nontrivial channel,
\cref{lem:root-profile-projectors} and subtraction of the trivial
rank-one summand give the projection onto
\(\Q(\delta_{x_\chi}-\delta_{y_\chi})\).  The cyclic rank-one lemma then
generates \(\End_\Q(V_\chi)\).  Project away these matrix blocks.  On the
remaining one-dimensional channels, the separating restrictions generate
all coordinate idempotents by interpolation.  Thus \(A\) contains each
simple factor of the displayed commutant.
\end{proof}

\begin{lemma}[Modular cyclicity]
\label{lem:modular-cyclicity}
Let \(R=\Z[1/M]\), let \(B\subseteq M_d(R)\) be generated by fixed
matrices, and let \(v\in R^d\).  Let \(\ell\nmid M\) be prime.  If words
\(w_1,\ldots,w_d\) in the generators satisfy
\[
 \det[w_1v\ \cdots\ w_dv]\not\equiv0\pmod\ell,
\]
then \(v\) is cyclic over \(\Q\).
\end{lemma}

\begin{proof}
The orbit determinant belongs to \(R\).  Its nonzero reduction cannot
come from the zero rational number.
\end{proof}

\subsection{The Monster symmetry commutant}

Let \(\Gamma_{\M}\) be the group of permutations of the \(194\) Monster
classes that preserve class size and commute with every prime-power map.  It
is the automorphism group of the marked weighted power system.

\begin{theorem}[Weighted-power symmetry and commutant]
\label{thm:monster-symmetry-commutant}
The group \(\Gamma_{\M}\) is elementary abelian of rank \(14\).  Its
subgroup generated by the nineteen stored ordinary-table automorphisms has
rank \(13\), and a complementary factor is generated by
the involution
\[
 (16B\ 16C)(32A\ 32B).
\]
For the permutation representation on
\(V_{\M}=\Q^{\operatorname{Conj}(\M)}\),
\[
 \operatorname{End}_{\Q\Gamma_{\M}}(V_{\M})
 =\Pram(\M)
 \cong
 M_{170}(\Q)\oplus M_5(\Q)\oplus M_3(\Q)
 \oplus M_2(\Q)^{\oplus4}\oplus\Q^{\oplus8}.
\]
Its dimension is \(28958\); its center and commutant have dimension \(15\).
\end{theorem}

\begin{proof}
Start with the coloring by class size and element order.  Refine it by the
colors of all prime-power images and by the weighted color counts in each
reverse fibre.  The first two refinement rounds have \(161\) and \(170\)
cells, and the second coloring is stable.  Its nonsingleton cells are
twenty-four pairs.  Commutation with the fifteen prime-power maps joins
their swap variables into fourteen components of sizes
\[
 2,5,3,2,1,2,2,1,1,1,1,1,1,1.
\]
Every component swap satisfies all power-map equations.  Hence
\(\Gamma_{\M}\cong(C_2)^{14}\).  The binary matrix of the nineteen stored
table automorphisms has rank \(13\); its missing coordinate is the component
\((16B\ 16C)(32A\ 32B)\).

Each singleton contributes one trivial vector.  Each pair contributes its
sum to the trivial channel and its difference to the character of its swap
component.  Thus
\[
 V_{\M}\cong
 \Q_{\mathbf1}^{170}\oplus
 \Q_{\chi_1}^{2}\oplus\Q_{\chi_2}^{5}\oplus
 \Q_{\chi_3}^{3}\oplus\Q_{\chi_4}^{2}\oplus
 \Q_{\chi_6}^{2}\oplus\Q_{\chi_7}^{2}
 \oplus\bigoplus_{i\in\{5,8,9,10,11,12,13,14\}}\Q_{\chi_i}.
\]
This gives the displayed target commutant.

It remains to prove saturation.  Since every Monster element has exponent
dividing \(N=\exp(\M)\), the operator
\(|\M|^{-1}P_N^*P_N\) is the coordinate projection onto \(1A\).
Modulo \(1000003\), the orbit of \(\delta_{1A}\) has rank \(170\), with a
retained basis of words of length at most six.  The pairs
\[
 16B/C,\ 23A/B,\ 31A/B,\ 39C/D,\ 44A/B,\ 47A/B
\]
have isolated root profiles for the six nontrivial channels of dimension
greater than one.  Their orbit determinants modulo \(1000003\) are
\[
 1,\ 2,\ 1,\ 1,\ -1,\ 1,
\]
with maximum word lengths \(1,3,1,1,1,1\).  Finally, the restrictions of
\((P_2,P_5,P_7)\) to the eight scalar channels are pairwise distinct.
The hypotheses of \cref{thm:root-cyclic-saturation} hold, so
\(\Pram(\M)=\End_{\Q\Gamma_{\M}}(V_{\M})\).  The certificate uses
\(170+2+5+3+2+2+2=186\) retained orbit vectors.  All denominators are
invertible modulo \(1000003\), so \cref{lem:modular-cyclicity} lifts the
seven orbit determinants to \(\Q\).
\end{proof}

\begin{remark}
The exceptional rational \(M_2(\Q)\)-channel is the unique symmetry
direction on which the complementary involution acts and the stored
table-automorphism subgroup acts trivially.  No identification of the
stored subgroup with the full unit-power image is used.
\end{remark}



\subsection{Symmetries invisible to the character table}
\label{sec:invisible-symmetries}

Galois conjugation of classes is an automorphism of the character table
and of the weighted power system, so \(\Psi(U_N)\subseteq\Gamma_G\) for
every \(G\), where \(\Gamma_G\) is the automorphism group of the weighted
power system.  \Cref{thm:monster-symmetry-commutant} shows that the
inclusion can be strict: for the Monster, \(\Gamma_{\M}\cong(C_2)^{14}\)
while the Galois image has rank \(13\), the thirteen distinct quadratic
characters of \cref{prop:monster-imaginary}.  The complementary involution
\((16B\ 16C)(32A\ 32B)\) preserves class sizes and all fifteen prime-power
maps.  It lies outside the group generated by the stored automorphisms of the
character table, so it is not induced by any automorphism of that table.
Thus the classes \(16B,16C\) and \(32A,32B\) are distinguished by the
character table but not by class sizes and power maps.

At the level of the algebra this involution is visible as the unique
nonprincipal block of \(\Pram(\M)\) on which the good labels act
trivially: by \cref{prop:rational-splitting} the Galois-fixed space
\(V_{\M}^{\Gamma}\) has dimension \(172\), and it decomposes under
\(\Pram(\M)\) as the \(170\)-dimensional principal block plus one
\(M_2(\Q)\)-block, the difference channel of the complementary involution.

The same signature occurs elsewhere in the census.  In the character table
of \cref{thm:good-label-census}, \(Co_2\) and \(Co_3\) each have an
\(M_2(\Q)\)-block carrying the trivial character, so each has a
two-dimensional nonprincipal simple factor inside its Galois-fixed space.
By \cref{thm:balanced-commutant} that factor comes from a balanced kernel
commuting with \(\Psi(U_N)\) and not accounted for by it.  Whether it
arises from a permutation of classes, as for the Monster, is not
determined by the computations reported here; the symmetry group
\(\Gamma_G\) was computed only for \(\M\).

For \(Fi_{23}\) the answer is known and is negative.  Its census row has a
scalar block with trivial character spanned by
\[
 v=\delta_{12H}-\tfrac13\,\delta_{12L}-\tfrac14\,\delta_{12M}
   +\tfrac1{12}\,\delta_{12O},
\]
a joint eigenvector of all sixteen generators, with eigenvalue \(0\) for
\(\Psi_2,\Psi_2^*,\Psi_3,\Psi_3^*\) and \(1\) for the others, orthogonal to
the principal block \(A\delta_{1A}\), which has exact dimension \(91\)
(\cref{sec:certification}).  The four coefficients of \(v\) have distinct
absolute values, so no permutation of classes commuting with the power maps
moves \(v\) except trivially; every element of \(\Q\Gamma_{Fi_{23}}\) acts
on the Galois-fixed space \(V^\Gamma\) as a scalar, while the projection
onto \(\Q v\) lies in the commutant and acts on \(V^\Gamma\) nontrivially.
Hence \(\Bal(Fi_{23})\supsetneq\Q\Gamma_{Fi_{23}}\): the commutant of
\(\Pram(Fi_{23})\) contains a balanced kernel that is not a combination of
permutation symmetries.  The Monster's extra block is a permutation
symmetry invisible to the character table; the \(Fi_{23}\) block is a
fractional symmetry invisible even to the automorphism group of the weighted
power system.

\subsection{Computational record}
\label{sec:characters-record}

The file \path{good_label_characters.json} records, per group, the generators
of \(U_N\) and the reconstruction data of \cref{thm:unit-reconstruction}, and
per primitive central idempotent the center field and its defining
polynomial, \(d_k\), the class support with its element orders, whether that
support is a single \(\Gamma\)-orbit, the values \(\lambda_k(u)\) at the
generators with their minimal polynomials and orders, and the conductor,
parity, and fundamental discriminant of \(\lambda_k\).  It checks
\eqref{eq:lefschetz}, \eqref{eq:trivial-block-sum},
\(\sigma_u\sigma_v=\sigma_{uv}\), and the admissible completions of
\cref{prop:ambiguity-harmless}; each generator has one completion.  Every
block acts by center scalars and every quadratic block agrees with its
Kronecker symbol at every generator.
